% GENERATED FILE. Edit canonical lessons, then run npm run generate:books.
% canonical-source-hash: fnv1a64:2ea4d423f437c979
% canonical-lessons: AR-C492-ghira, AR-C492-ghayur, AR-C492-hasud, AR-C492-hasada-envy, AR-C492-ghara

\chapter{Envy and jealousy}
\label{ch:ar-envy-and-jealousy}
\hlchaptermodality{492}

\begin{chapteropening}
\textbf{By the end of this chapter:} \emph{I can talk about jealousy and envy, and describe a jealous or envious person.}

Complete the last lesson of chapter 492: I can talk about jealousy and envy, and describe a jealous or envious person.
\end{chapteropening}

\section[\texorpdfstring{\ar{غيرة} (ghīra)}{غيرة (ghīra)}]{\ar{غيرة} (ghīra) — jealousy}

\label{lesson:AR-C492-ghira}



\begin{quote}
Before the new one: say the Arabic for grateful, then the Arabic for to impress.
\end{quote}

\subsection*{You'll want to know: \ar{غيرة}}
\textbf{\ar{غيرة}} — \emph{ghīra} — \textquotedblleft{}jealousy\textquotedblright{}.

\begin{quote}
\ar{شعر} \ar{الطفل} \ar{بالغيرة} \ar{من} \ar{أخيه} \ar{الصغير}.
\end{quote}

\begin{quote}
\emph{shaʿara aṭ-ṭifl bi-l-ghīra min akhīhi aṣ-ṣaghīr.}
\end{quote}

\begin{quote}
The child felt jealous of his little brother.
\end{quote}

\begin{cousinweb}
\textbf{\ar{غيرة}} (\emph{ghīra}) is built on the native Arabic root \emph{gh-y-r}, \textquotedblleft{}jealousy, a jealous care for one's own\textquotedblright{} (Lane). Ghīra is jealousy, often a protective one.
\end{cousinweb}

\begin{grammarlens}[title={What you've built}]
One more word for describing what you lived and felt.
\end{grammarlens}

\subsection*{Your turn}
\begin{itemize}
\raggedright
  \item \emph{Say it:} \emph{ghīra}
  \item \emph{Say it:} \emph{ghīra}, once more
  \item \emph{Say it:} \emph{shaʿara aṭ-ṭifl bi-l-ghīra min akhīhi aṣ-ṣaghīr.} — The child felt jealous of his little brother.
  \item \emph{Recall:} say the Arabic for grateful, then the Arabic for to impress, then say \emph{ghīra} again
\end{itemize}

\begin{tcolorbox}[breakable,colback=teal!4,colframe=teal!35!black,title={Before you move on}]
What does \ar{غيرة} mean? (\textquotedblleft{}Jealousy\textquotedblright{}.) Say it once more.
\end{tcolorbox}
\section[\texorpdfstring{\ar{غيور} (ghayūr)}{غيور (ghayūr)}]{\ar{غيور} (ghayūr) — jealous}

\label{lesson:AR-C492-ghayur}



\begin{quote}
Before the new one: say the Arabic for to impress, then the Arabic for jealousy.
\end{quote}

\subsection*{You'll want to know: \ar{غيور}}
\textbf{\ar{غيور}} — \emph{ghayūr} — \textquotedblleft{}jealous\textquotedblright{}.

\begin{quote}
\ar{زوجها} \ar{رجل} \ar{غيور}.
\end{quote}

\begin{quote}
\emph{zawjuhā rajul ghayūr.}
\end{quote}

\begin{quote}
Her husband is a jealous man.
\end{quote}

\begin{cousinweb}
\textbf{\ar{غيور}} (\emph{ghayūr}) is built on the native Arabic root \emph{gh-y-r}, \textquotedblleft{}jealousy, a jealous care for one's own\textquotedblright{} (Lane). Ghayūr is the pattern of a lasting trait.
\end{cousinweb}

\begin{grammarlens}[title={What you've built}]
One more word for describing what you lived and felt.
\end{grammarlens}

\subsection*{Your turn}
\begin{itemize}
\raggedright
  \item \emph{Say it:} \emph{ghayūr}
  \item \emph{Say it:} \emph{ghayūr}, once more
  \item \emph{Say it:} \emph{zawjuhā rajul ghayūr.} — Her husband is a jealous man.
  \item \emph{Recall:} say the Arabic for to impress, then the Arabic for jealousy, then say \emph{ghayūr} again
\end{itemize}

\begin{tcolorbox}[breakable,colback=teal!4,colframe=teal!35!black,title={Before you move on}]
What does \ar{غيور} mean? (\textquotedblleft{}Jealous\textquotedblright{}.) Say it once more.
\end{tcolorbox}
\section[\texorpdfstring{\ar{حسود} (ḥasūd)}{حسود (ḥasūd)}]{\ar{حسود} (ḥasūd) — envious}

\label{lesson:AR-C492-hasud}



\begin{quote}
Before the new one: say the Arabic for jealousy, then the Arabic for jealous.
\end{quote}

\subsection*{You'll want to know: \ar{حسود}}
\textbf{\ar{حسود}} — \emph{ḥasūd} — \textquotedblleft{}envious\textquotedblright{}.

\begin{quote}
\ar{لا} \ar{أحب} \ar{الجار} \ar{الحسود}.
\end{quote}

\begin{quote}
\emph{lā uḥibbu al-jār al-ḥasūd.}
\end{quote}

\begin{quote}
I do not like the envious neighbour.
\end{quote}

\begin{cousinweb}
\textbf{\ar{حسود}} (\emph{ḥasūd}) is built on the native Arabic root \emph{ḥ-s-d}, \textquotedblleft{}to envy\textquotedblright{} (Lane). Ḥasūd is the pattern of a lasting trait: an envious person.
\end{cousinweb}

\begin{grammarlens}[title={What you've built}]
One more word for describing what you lived and felt.
\end{grammarlens}

\subsection*{Your turn}
\begin{itemize}
\raggedright
  \item \emph{Say it:} \emph{ḥasūd}
  \item \emph{Say it:} \emph{ḥasūd}, once more
  \item \emph{Say it:} \emph{lā uḥibbu al-jār al-ḥasūd.} — I do not like the envious neighbour.
  \item \emph{Recall:} say the Arabic for jealousy, then the Arabic for jealous, then say \emph{ḥasūd} again
\end{itemize}

\begin{tcolorbox}[breakable,colback=teal!4,colframe=teal!35!black,title={Before you move on}]
What does \ar{حسود} mean? (\textquotedblleft{}Envious\textquotedblright{}.) Say it once more.
\end{tcolorbox}
\section[\texorpdfstring{\ar{حسد} (ḥasada)}{حسد (ḥasada)}]{\ar{حسد} (ḥasada) — to envy}

\label{lesson:AR-C492-hasada-envy}



\begin{quote}
Before the new one: say the Arabic for jealous, then the Arabic for envious.
\end{quote}

\subsection*{You'll want to know: \ar{حسد}}
\textbf{\ar{حسد}} — \emph{ḥasada} — \textquotedblleft{}to envy\textquotedblright{}.

\begin{quote}
\ar{حسد} \ar{الموظف} \ar{زميله} \ar{على} \ar{راتبه}.
\end{quote}

\begin{quote}
\emph{ḥasada al-muwaẓẓaf zamīlahu ʿalā rātibihi.}
\end{quote}

\begin{quote}
The employee envied his colleague his salary.
\end{quote}

Present: \ar{يحسد}, \emph{yaḥsud}.

It takes \ar{على} for the thing envied.

\begin{cousinweb}
\textbf{\ar{حسد}} (\emph{ḥasada}) is built on the native Arabic root \emph{ḥ-s-d}, \textquotedblleft{}to envy\textquotedblright{} (Lane). Lane gives yaḥsudu as the usual present.
\end{cousinweb}

\begin{grammarlens}[title={What you've built}]
One more word for describing what you lived and felt.
\end{grammarlens}

\subsection*{Your turn}
\begin{itemize}
\raggedright
  \item \emph{Say it:} \emph{ḥasada}
  \item \emph{Say it:} \emph{ḥasada}, once more
  \item \emph{Say it:} \emph{ḥasada al-muwaẓẓaf zamīlahu ʿalā rātibihi.} — The employee envied his colleague his salary.
  \item \emph{Recall:} say the Arabic for jealous, then the Arabic for envious, then say \emph{ḥasada} again
\end{itemize}

\begin{tcolorbox}[breakable,colback=teal!4,colframe=teal!35!black,title={Before you move on}]
What does \ar{حسد} mean? (\textquotedblleft{}To envy\textquotedblright{}.) Say it once more.
\end{tcolorbox}
\section[\texorpdfstring{\ar{غار} (ghāra)}{غار (ghāra)}]{\ar{غار} (ghāra) — to be jealous}

\label{lesson:AR-C492-ghara}



\begin{quote}
Before the new one: say the Arabic for envious, then the Arabic for to envy.
\end{quote}

\subsection*{You'll want to know: \ar{غار}}
\textbf{\ar{غار}} — \emph{ghāra} — \textquotedblleft{}to be jealous\textquotedblright{}.

\begin{quote}
\ar{تغار} \ar{أختي} \ar{عندما} \ar{يلعب} \ar{أبي} \ar{مع} \ar{أخي}.
\end{quote}

\begin{quote}
\emph{taghāru ukhtī ʿindamā yalʿabu abī maʿa akhī.}
\end{quote}

\begin{quote}
My sister gets jealous when my father plays with my brother.
\end{quote}

Present: \ar{يغار}, \emph{yaghār}.

It takes \ar{من} or \ar{على}: \ar{غار} \ar{منه}. The same letters also spell \ar{غار}, ghār, a cave, from another root.

\begin{cousinweb}
\textbf{\ar{غار}} (\emph{ghāra}) is built on the native Arabic root \emph{gh-y-r}, \textquotedblleft{}jealousy, a jealous care for one's own\textquotedblright{} (Lane). It is the verb that goes with ghīra and ghayūr.
\end{cousinweb}

\begin{grammarlens}[title={What you've built}]
One more word for describing what you lived and felt.
\end{grammarlens}

\subsection*{Your turn}
\begin{itemize}
\raggedright
  \item \emph{Say it:} \emph{ghāra}
  \item \emph{Say it:} \emph{ghāra}, once more
  \item \emph{Say it:} \emph{taghāru ukhtī ʿindamā yalʿabu abī maʿa akhī.} — My sister gets jealous when my father plays with my brother.
  \item \emph{Recall:} say the Arabic for envious, then the Arabic for to envy, then say \emph{ghāra} again
\end{itemize}

\begin{tcolorbox}[breakable,colback=teal!4,colframe=teal!35!black,title={Before you move on}]
What does \ar{غار} mean? (\textquotedblleft{}To be jealous\textquotedblright{}.) Say it once more.
\end{tcolorbox}
